\documentclass{article}
\usepackage[margin=1.5cm]{geometry}
\usepackage{amsmath}
\usepackage{enumitem}

\setlist{nosep,leftmargin=*}

\begin{document}
\twocolumn

\title{Chapter 6 Warm-Up: Functions and a Sawtooth Wave}
\author{Prof. Jordan C. Hanson}
\date{}

\maketitle

\section{Defining a Function}

\begin{enumerate}

\item Python functions allow us to give a name to a calculation and
evaluate it for different arguments.  Begin with

\begin{verbatim}
import numpy as np
import matplotlib.pyplot as plt

def wave(x):
    return np.sin(x)
\end{verbatim}

Create 400 values between $-\pi$ and $\pi$:

\begin{verbatim}
x = np.linspace(-np.pi,
                np.pi, 400)

y = wave(x)
\end{verbatim}

Here, the argument \texttt{x} is an array, so the function evaluates
$\sin(x)$ at every value in the array.

\end{enumerate}


\section{Adding Trigonometric Terms}

\begin{enumerate}

\item The periodic sawtooth function

\[
f(x)=x,\qquad -\pi < x < \pi
\]

can be approximated by the sum

\[
f_N(x)=
2\sum_{n=1}^{N}
\frac{(-1)^{n+1}}{n}\sin(nx).
\]

Define a Python function that calculates this sum:

\begin{verbatim}
def sawtooth(x, N):
    total = 0.0

    for n in range(1, N + 1):
        term = 2 * (-1)**(n + 1)
        term = term * np.sin(n*x)/n
        total = total + term

    return total
\end{verbatim}

\begin{itemize}

\item What are the two parameters of this function?

\item What does \texttt{n} represent?

\item Why must \texttt{total} be defined before the loop?

\item What value is returned by the function?

\end{itemize}

\end{enumerate}


\section{Code Lab: Build the Sawtooth}

\begin{enumerate}

\item Evaluate the approximation with only one term:

\begin{verbatim}
y1 = sawtooth(x, 1)
\end{verbatim}

Now evaluate it using more terms:

\begin{verbatim}
y3 = sawtooth(x, 3)
y10 = sawtooth(x, 10)
\end{verbatim}

Plot all three results:

\begin{verbatim}
plt.plot(x, y1, label="N = 1")
plt.plot(x, y3, label="N = 3")
plt.plot(x, y10, label="N = 10")

plt.xlabel("x")
plt.ylabel("f(x)")
plt.legend()
plt.grid()
plt.show()
\end{verbatim}

\begin{itemize}

\item Which approximation looks most like a straight line between
$-\pi$ and $\pi$?

\item What happens near $x=-\pi$ and $x=\pi$?

\item Increase the number of terms to

\begin{verbatim}
N = 50
\end{verbatim}

and plot the result.  Does adding terms make the transition at the
edge sharper?

\end{itemize}

\end{enumerate}


\section{Challenge}

\begin{enumerate}

\item Modify the plotting code so that it compares

\[
N=1,\qquad N=5,\qquad N=20.
\]

Then add the exact line $f(x)=x$:

\begin{verbatim}
plt.plot(x, x, label="f(x) = x")
\end{verbatim}

Which value of $N$ gives the closest approximation over most of the
interval?

\end{enumerate}

\end{document}